\chapter{Triển khai, tích hợp và minh chứng kết quả}
\label{ch:trien-khai}

\section{Tích hợp hệ thống}

\begin{sloppypar}

\begin{itemize}
  \item \textbf{Dùng trực tiếp}: \texttt{VectorStoreIndex}, \texttt{StorageContext}, \texttt{SentenceSplitter},
        \texttt{QueryFusionRetriever}, \texttt{SentenceTransformerRerank}, \texttt{SimilarityPostprocessor},
        \texttt{RetrieverQueryEngine}, \texttt{RetrieverEvaluator}, các evaluator câu trả lời.
  \item \textbf{Viết lại / tự viết}: bước dữ liệu (OCR + Markdown hai tầng thay \texttt{SimpleDirectoryReader}),
        loader Node theo heading, BM25, client model server, độ đo thống kê.
  \item \textbf{Điểm vào}: \texttt{src/pipeline.py} (dòng lệnh), \texttt{app/app.py} (Gradio, ghi nhật ký câu hỏi),
        \texttt{eval/run\_batch\_cau\_hoi.py} (chạy hàng loạt).
  \item \textbf{Thí nghiệm chỉ để ra quyết định}: so sánh vector/rerank/hybrid (quyết định tắt cả hai), chọn
        6 văn bản đang áp dụng thay vì nạp cả văn bản hết hiệu lực (văn bản cũ lấn át Quy chế 2021 khi truy
        hồi), chọn profile \texttt{server} làm cấu hình đo chính.
\end{itemize}
\end{sloppypar}

\section{Kết quả đạt được}

\subsection{Truy hồi}

Bảng~\ref{tab:recall-vector} và~\ref{tab:baseline} (Chương~\ref{ch:cong-nghe-cot-loi}) cho chỉ số truy hồi;
Bảng~\ref{tab:top-k} cho top-$k$ Node của các câu mẫu; Bảng~\ref{tab:recall} (Chương~\ref{ch:nang-cao}) so
sánh các cấu hình nâng cao.

\subsection{Câu trả lời}
\label{sec:danh-gia-tra-loi}

\subsubsection{Cách đánh giá}

Bản báo cáo trước chấm câu trả lời bằng tay trên một phiên hỏi đáp 56 câu. Bản này bổ sung
\texttt{eval/evaluate\_answers.py} để chấm tự động bằng chính các evaluator của LlamaIndex, chạy hàng
loạt qua \texttt{BatchEvalRunner}:

\begin{itemize}
  \item \texttt{FaithfulnessEvaluator} --- câu trả lời có suy ra được từ ngữ cảnh truy hồi không;
  \item \texttt{RelevancyEvaluator} --- câu trả lời có đúng câu hỏi không;
  \item \texttt{CorrectnessEvaluator} --- so với đáp án chuẩn trong \texttt{eval/cau\_hoi\_danh\_gia.json};
  \item chế độ \texttt{closed-book} (LLM trả lời \emph{không} có ngữ cảnh) làm mốc so sánh RAG với
        không RAG.
\end{itemize}

Bộ câu hỏi chấm câu trả lời gồm 13 câu, mỗi câu có trường \texttt{pham\_vi} (``trong'' / ``ngoai'')
và đáp án chuẩn lấy từ bảng đã kiểm chứng thủ công ở bản báo cáo trước. Nhờ trường \texttt{pham\_vi},
script tách \textbf{từ chối đúng} (câu ngoài phạm vi) khỏi \textbf{từ chối sai} (đáp án có trong ngữ
cảnh nhưng hệ thống vẫn nói không tìm thấy) --- chính là hai loại đã bị gộp thành một con số ``52\%
từ chối'' ở bản trước.

\paragraph{Trạng thái chạy.} Script đã viết xong, kiểm thử import và logic nhận dạng từ chối, nhưng
\textbf{chưa chạy} \texttt{make eval-answers}. Trong tuần 05--09/10, LLM đã chạy được qua Ollama trên
Colab (\texttt{scripts/setup\_server\_v2.sh}), và nhóm dùng nó để chạy lại 56 câu hỏi của phiên trước
với prompt mới (Mục~\ref{sec:chay-lai-56-cau}) --- vẫn chấm tay. Vì vậy Mục~\ref{sec:tra-loi-thu-cong}
trình bày kết quả chấm tay của phiên chạy trước (có ghi rõ nguồn), còn các bảng điểm
faithfulness/relevancy/correctness sẽ được điền khi chạy \texttt{make eval-answers}.

\subsubsection{Kết quả chấm tay (phiên chạy 56 câu, profile \texttt{server} có rerank)}
\label{sec:tra-loi-thu-cong}

Nhật ký \texttt{eval/results/test-basic.md} ghi một phiên chạy \texttt{src/pipeline.py}: 56 câu hỏi do
nhóm tự đặt theo cách sinh viên thật sự hỏi, gồm cả viết tắt (``bao nhiêu chỉ''), sai chính tả
(``xuất sắt''), câu hỏi nối tiếp (``còn kỹ sư'') và câu ngoài phạm vi. Hệ thống trả lời ``không tìm
thấy'' ở 29/56 câu (52\%) và đưa ra câu trả lời ở 27 câu.

\begin{table}[H]
  \centering
  \small
  \begin{tabularx}{\textwidth}{>{\hsize=0.95\hsize}X >{\hsize=0.95\hsize}X >{\hsize=1.1\hsize}X l}
    \toprule
    \textbf{Câu hỏi} & \textbf{Hệ thống trả lời} & \textbf{Quy chế 2021 quy định} & \textbf{Đánh giá} \\
    \midrule
    Học cử nhân cần bao nhiêu tín chỉ & Tối thiểu 120 tín chỉ & 120 tín chỉ (Điều 2) & Đúng \\
    Muốn được B cần bao nhiêu điểm & Từ 7,0 đến 8,4 & B: 7,0--8,4 (Điều 9) & Đúng \\
    Cảnh báo kỳ 1 và 3, kỳ 2 bình thường có bị buộc thôi học & Không, vì hai lần không liên tiếp & Buộc thôi học khi quá hai lần cảnh báo \emph{liên tiếp} (Điều 11) & Đúng \\
    Mua bằng tiếng Anh giả nộp cho trường & Không được; buộc thôi học, thu hồi bằng & Như câu trả lời (Điều 19) & Đúng \\
    \midrule
    8,4 hệ 10 thì hệ 4 được bao nhiêu & 3,4 điểm & Khóa trước 2025: 8,4 là B, quy đổi 3,0; khóa 2025 trở đi (QĐ 3183): 8,4 là B+, quy đổi 3,5 (Điều 9, 10) & Sai \\
    Khối lượng học kỳ chính & Tối thiểu 30 tuần & 2 học kỳ chính \emph{cộng lại} tối thiểu 30 tuần (Điều 6) & Sai \\
    Học bị F thì làm gì & Học lại; điểm lần học cuối là điểm chính thức & Học phần không đạt phải học lại, điểm \emph{lần học cuối} là điểm chính thức; ``cao nhất'' chỉ áp dụng khi học cải thiện (Điều 9 khoản 4) & Đúng$^{*}$ \\
    Học điểm P thì sao & P không tính vào điểm TB, ``quy đổi thành 0'' & P không được tính vào điểm TB (Điều 10 khoản 3), không quy đổi & Sai một phần \\
    \midrule
    Ra trường xuất sắc cần bao nhiêu điểm & Không tìm thấy & Từ 3,6 đến 4,0 (Điều 10), và Điều 10 \emph{có} trong 3 Điều nguồn & Từ chối sai \\
    Còn kỹ sư (hỏi tiếp câu cử nhân) & Không tìm thấy & 150 tín chỉ; hỏi lại đầy đủ thì trả lời đúng & Từ chối sai \\
    \midrule
    Điểm rèn luyện là gì & Không tìm thấy & Không thuộc Quy chế đào tạo (văn bản riêng của CTSV) & Từ chối đúng \\
    Bao nhiêu tiền một tín chỉ & Không tìm thấy & Không thuộc Quy chế đào tạo & Từ chối đúng \\
    \bottomrule
  \end{tabularx}
  \caption{Một số câu trả lời trong nhật ký phiên chạy trước, đối chiếu với văn bản Quy chế 2021.
  Đây là số liệu chấm tay, không sinh lại được bằng lệnh. $^{*}$Bản báo cáo trước chấm câu này là
  ``Sai một phần''; khi đối chiếu lại khoản 4 Điều 9 trong tuần 05--09/10 thì câu trả lời của hệ thống
  là đúng, còn đáp án tham chiếu đã nhầm hai trường hợp học lại và học cải thiện.}
  \label{tab:tra-loi}
\end{table}

\subsubsection{Phân tích}

\paragraph{Prompt chống bịa hoạt động, nhưng hơi quá tay.} Với câu hỏi ngoài phạm vi, hệ thống từ chối
thay vì bịa --- đúng mục tiêu. Nhưng tỷ lệ từ chối 52\% cao hơn tỷ lệ câu thật sự ngoài phạm vi: có
câu mà thông tin \emph{đã nằm trong ngữ cảnh} vẫn bị từ chối, như câu hỏi về điểm ra trường loại xuất
sắc (viết sai chính tả ``xuất sắt''). Đây chính là lý do phải tách ``từ chối đúng'' khỏi ``từ chối
sai'' như thiết kế ở Mục~\ref{sec:danh-gia-tra-loi}.

\paragraph{Lỗi sai tập trung ở phép tính và chi tiết số.} Các câu trả lời sai đều có Điều đúng trong
ngữ cảnh (truy hồi đúng), nhưng LLM đọc sai hoặc suy diễn thêm: tự ``quy đổi'' 8,4 thành 3,4 thay vì
tra bảng, hay đọc ``30 tuần cho 2 học kỳ'' thành ``30 tuần mỗi học kỳ''. Đây là loại lỗi nguy hiểm nhất vì câu
trả lời vẫn trích dẫn đúng Điều. Bản này bổ sung hai ràng buộc trong \texttt{QA\_PROMPT} nhắm đúng
loại lỗi này: \emph{trích nguyên văn câu quy định trước khi kết luận} và \emph{không tự quy đổi thang
điểm, cần tra bảng thì chép đúng hàng của bảng}. Hiệu quả của hai ràng buộc này cần được đo bằng
\texttt{make eval-answers} trên máy có LLM, và được ghi nhận là việc còn lại.

\paragraph{Câu trả lời không ổn định --- đã xử lý ở mức mã nguồn.} Bản trước ghi nhận cùng một câu hỏi
chỉ khác một dấu hỏi chấm mà lần đầu từ chối, lần sau trả lời ``Không''. Nguyên nhân là
\texttt{Ollama(...)} không đặt \texttt{temperature}, nên mô hình lấy mẫu ngẫu nhiên. Bản này đặt
\texttt{temperature=0} và \texttt{seed=42} trong \texttt{make\_llm()}; đây là sửa đổi đã kiểm tra ở
mức mã nguồn nhưng cần một phiên chạy LLM mới để xác nhận bằng số.

\paragraph{Không có ngữ cảnh hội thoại.} \texttt{RetrieverQueryEngine} gửi từng câu hỏi độc lập, nên
câu nối tiếp ``còn kỹ sư'' bị từ chối. Sinh viên thường hỏi nối tiếp, nên đây vẫn là hạn chế lớn nhất
ở tầng sinh.

\subsubsection{Chạy lại 56 câu với prompt mới (tuần 05--09/10)}
\label{sec:chay-lai-56-cau}

\paragraph{Cách chạy.} Cùng 56 câu hỏi của phiên trước (\texttt{eval/results/cau-hoi01.md}) được chạy
lại bằng \texttt{eval/run\_batch\_cau\_hoi.py} (bản notebook: \texttt{notebooks/chay\_cau\_hoi01.ipynb}),
profile \texttt{server}, \texttt{qwen3:8b} trên Colab, \texttt{temperature=0}, \texttt{seed=42}. Kết quả
ghi ở \texttt{eval/results/cau-hoi01-result.md} (phiên bản trong commit \texttt{be7b81e}, ngày 08/10).
Prompt dùng trong lần chạy này đã có hai ràng buộc ``trích nguyên văn'' và ``không tự quy đổi'' nhưng
\emph{chưa} có chú giải thang điểm ở Mục~\ref{sec:nham-thang-diem}.

\paragraph{Lưu ý khi so sánh.} Hai phiên khác nhau ở \emph{hai} điểm cùng lúc: prompt mới, và cấu hình
ngữ cảnh --- phiên trước bật rerank (LLM chỉ nhận 3 Node), phiên này tắt rerank (LLM nhận 10 Node).
Vì vậy chênh lệch dưới đây không quy được riêng cho prompt; muốn tách phải chạy thêm một phiên chỉ đổi
một yếu tố.

\paragraph{Kết quả tổng.} Số câu trả lời ``không tìm thấy'' giảm từ 29/56 (52\%) xuống 23/56 (41\%).

\begin{table}[H]
  \centering
  \small
  \begin{tabularx}{\textwidth}{>{\hsize=0.9\hsize}X >{\hsize=1.1\hsize}X >{\hsize=1.1\hsize}X >{\hsize=0.9\hsize}X}
    \toprule
    \textbf{Câu hỏi} & \textbf{Trả lời (phiên 08/10)} & \textbf{Quy chế quy định} & \textbf{So với phiên trước} \\
    \midrule
    Ra trường xuất sắt cần bao nhiêu điểm & ĐTB tích lũy 3,6--4,0 & 3,6--4,0 (Điều 10 khoản 4) & Từ chối sai $\to$ \textbf{Đúng} \\
    8,4 hệ 10 thì hệ 4 được bao nhiêu & 3,5 & 3,5 với khóa 2025 trở đi (B+); 3,0 với khóa trước & Sai $\to$ Đúng một phần (không nêu khóa) \\
    Điểm D+ ứng với khoảng nào trên thang 10 & 4,8--5,4 & 4,8--5,4 (QĐ 3183, Điều 1 khoản 3) & (câu mới) Đúng \\
    Giỏi tiếng Anh nhưng không có chứng chỉ có ra trường được không & Được xét nếu đủ điều kiện; chưa đạt chuẩn ngoại ngữ thì có 3 năm để hoàn thiện & Như câu trả lời (Điều 13 khoản 4) & Đúng \\
    \midrule
    Không có chứng chỉ tiếng Anh có ra trường được không & Không, dẫn Điều 3 về học chương trình thứ hai ngành Ngôn ngữ Anh & Điều 13 khoản 4, như câu trên & \textbf{Sai} (lấy nhầm văn bản QĐ 810) \\
    Học bị F thì làm gì & F quy đổi 0 và ``không được tính vào điểm trung bình'' & F quy đổi 0 và \emph{được} tính (Điều 10 khoản 2--3) & Đúng $\to$ \textbf{Sai} \\
    Khối lượng học kỳ chính & Mỗi học kỳ chính tối thiểu 30 tuần & 2 học kỳ chính cộng lại tối thiểu 30 tuần (Điều 6) & Sai (không đổi) \\
    Rớt một môn, các kỳ còn lại 4.0 có được xuất sắc không & Không tìm thấy & Chỉ bị giảm hạng nếu học lại vượt 5\% tổng tín chỉ (Điều 13 khoản 3) & Từ chối sai \\
    Còn kỹ sư (hỏi tiếp) & Không tìm thấy & 150 tín chỉ & Từ chối sai (không đổi) \\
    \midrule
    Mua bằng tiếng Anh giả (hai cách hỏi) & Một cách hỏi: đúng, dẫn Điều 19 khoản 3; cách kia: không tìm thấy & Buộc thôi học, thu hồi bằng (Điều 19 khoản 3) & Không ổn định theo cách diễn đạt \\
    Bao nhiêu tiền 01 tín & In ra nguyên chuỗi \texttt{/no\_think} & Ngoài phạm vi & Từ chối đúng $\to$ \textbf{Lỗi} \\
    \bottomrule
  \end{tabularx}
  \caption{Một số câu của phiên chạy lại ngày 08/10, chấm tay theo văn bản trong \texttt{data/processed}.}
  \label{tab:chay-lai-56-cau}
\end{table}

\paragraph{Phân tích.}
\begin{itemize}
  \item \textbf{Ít từ chối sai hơn}, rõ nhất ở câu ``xuất sắt'' --- trước đây bị từ chối dù Điều 10 có
        trong ngữ cảnh. Do ngữ cảnh tăng từ 3 lên 10 Node cùng lúc với đổi prompt, chưa kết luận được
        yếu tố nào có tác dụng.
  \item \textbf{Phạm vi hiệu lực theo khóa tuyển sinh là nguồn sai mới.} QĐ 3183 chỉ áp dụng thang điểm
        có B+, C+, D+ cho khóa tuyển sinh từ 2025 (Điều 2 của QĐ). Hệ thống trả lời theo một trong hai
        bảng mà không nói bảng đó áp dụng cho khóa nào: ``8,4 $\to$ 3,5'' đúng với khóa 2025, còn ``B là
        7,0--8,4'' chỉ đúng với khóa trước 2025. Prompt hiện chỉ nói ``có văn bản sửa đổi thì dùng nội
        dung đã sửa đổi'', tức là ngầm giả định mọi sinh viên đều thuộc khóa mới.
  \item \textbf{Ngữ cảnh rộng hơn cũng kéo theo văn bản nhiễu.} Câu ``không có chứng chỉ tiếng Anh'' lấy
        Điều 3 của QĐ 810 (điều kiện học chương trình thứ hai ngành Ngôn ngữ Anh) làm căn cứ, trong khi
        Điều 13 --- căn cứ đúng --- cũng có trong ngữ cảnh. Câu về điểm F sai ở chi tiết ``không tính vào
        điểm trung bình'', nhầm F với các điểm chữ không quy đổi ở khoản 3 Điều 10.
  \item \textbf{Không ổn định theo cách diễn đạt.} \texttt{temperature=0} làm \emph{cùng một} câu hỏi cho
        cùng câu trả lời, nhưng hai cách hỏi gần nhau vẫn có thể cho một câu đúng và một câu ``không tìm
        thấy''. Tính tất định không thay được độ vững theo diễn đạt.
  \item \textbf{Rò chỉ thị \texttt{/no\_think}.} Ở một câu ngoài phạm vi, Qwen3 trả về đúng chuỗi điều khiển
        \texttt{/no\_think} thay cho câu trả lời. Cần lọc chuỗi này ở đầu ra và đếm nó như một lỗi trong
        \texttt{evaluate\_answers.py}.
\end{itemize}


\subsection{So sánh trước -- sau các tối ưu}

\begin{table}[H]
  \centering
  \small
  \begin{tabularx}{\textwidth}{>{\hsize=0.8\hsize}X >{\hsize=1.1\hsize}X >{\hsize=1.1\hsize}X}
    \toprule
    \textbf{Vấn đề} & \textbf{Trước} & \textbf{Sau} \\
    \midrule
    Điều vắt qua hai trang & Bị chia đôi (mỗi trang một Document) & Một Node mỗi Điều (Markdown hai tầng) \\
    Điều 1--3 của quyết định ban hành & Trùng nhãn với Điều 1--3 của quy chế & Tách \texttt{phan="ban\_hanh"} \\
    Metadata rỗng trong embedding & Có (\texttt{chuong:}, \texttt{dieu:} rỗng) & Chỉ \texttt{nguon\_trich} \\
    Dùng nhầm chỉ mục cũ & Có thể, nếu quên \texttt{--rebuild} & Fingerprint tự build lại \\
    Câu trả lời giữa các lần hỏi & Không ổn định & \texttt{temperature=0}, \texttt{seed=42} \\
    Embedding + reranker trên GPU 4\,GB & \texttt{CUDA out of memory} & fp16, khoảng 2,2\,GB \\
    Tỷ lệ ``không tìm thấy'' (56 câu) & 52\% & 41\% (đổi prompt và top-$k$ cùng lúc) \\
    \bottomrule
  \end{tabularx}
  \caption{Các thay đổi đã đo hoặc kiểm chứng được.}
\end{table}

\subsection{Hướng dẫn chạy lại}

\begin{lstlisting}[caption={Các lệnh chạy lại (từ thư mục gốc repo).}, label={lst:chay-lai}]
python3 -m venv venv && venv/bin/pip install -r requirements.txt
bash scripts/pull_models.sh server            # hoặc bash scripts/setup_server_v2.sh trên Colab
RAG_PROFILE=server venv/bin/python src/index_builder.py
RAG_PROFILE=server venv/bin/python src/pipeline.py      # hỏi đáp dòng lệnh
RAG_PROFILE=server venv/bin/python app/app.py           # giao diện Gradio
make test && make eval && make baseline                 # kiểm thử và số liệu truy hồi
\end{lstlisting}

\section{Bảng phủ kỹ thuật}
\label{sec:bang-phu}

Bảng phản ánh trạng thái tuần 5; bản cuối sẽ khóa cùng commit niêm phong, nộp ít nhất 24 giờ trước buổi bảo vệ. Chỉ khai
những gì có trong commit đó.

\begin{table}[H]
  \centering
  \scriptsize
  \begin{tabularx}{\textwidth}{c >{\hsize=1.1\hsize}X c >{\hsize=0.7\hsize}X >{\hsize=1.2\hsize}X >{\hsize=1\hsize}X l}
    \toprule
    \textbf{TT} & \textbf{Nội dung} & \textbf{Tầng} & \textbf{Chức năng} & \textbf{File} & \textbf{Minh chứng} & \textbf{Trạng thái} \\
    \midrule
    1 & OCR + Markdown hai tầng & 1A & Nạp quy chế & \texttt{scripts/extract\_to\_md.py}, \texttt{clean\_md.py} & 6 văn bản & Dùng \\
    2 & Node theo Điều & 1A+1B & Đơn vị truy hồi & \texttt{src/processed\_loader.py} & 54 Node, kiểm thử & Dùng \\
    3 & Embedding + index + persist & 1A+1B & Lập chỉ mục & \texttt{src/index\_builder.py} & 77\,s / 1,5\,s & Dùng \\
    4 & Retriever top-$k$ & 1A & Truy hồi & \texttt{src/retriever.py} & Bảng~\ref{tab:top-k} & Dùng \\
    5 & Query engine + trích nguồn & 1A+1B & Hỏi đáp & \texttt{src/query\_engine.py} & Dòng ``Nguồn'' & Dùng \\
    6 & Recall@$k$, MRR, nDCG, CI & 1A & Kiểm chứng & \texttt{eval/evaluate.py}, \texttt{metrics.py} & R@1 81\%, R@3 100\% & Dùng \\
    7 & Từ chối ngoài corpus & 1B & Hỏi đáp & \texttt{src/query\_engine.py} & 41\% ``không tìm thấy'' & Chưa đo tỷ lệ đúng \\
    8 & Hybrid BM25 & 3 & Truy hồi & \texttt{src/bm25.py} & Bảng~\ref{tab:recall} & Có, tắt \\
    9 & Rerank cross-encoder & 3 & Truy hồi & \texttt{src/retriever.py} & Bảng~\ref{tab:recall} & Có, tắt \\
    10 & Model server HTTP & 3 & Hạ tầng & \texttt{src/model\_server.py} & Chạy trên Colab & Dùng \\
    11 & \texttt{.env}, 2 profile & 2 & Cấu hình & \texttt{src/config.py}, \texttt{.env.example} & --- & Dùng \\
    12 & Unit test + CI & 2 & Kiểm thử & \texttt{tests/}, \texttt{.github/workflows/ci.yml} & 26 kiểm thử & Dùng \\
    13 & Giao diện Gradio & 1B & Demo & \texttt{app/app.py} & Nhật ký câu hỏi & Chạy thử \\
    \bottomrule
  \end{tabularx}
  \caption{Bảng phủ kỹ thuật.}
  \label{tab:bang-phu}
\end{table}
